\documentclass[11pt,letterpaper]{article}
\usepackage[margin=0.7in,headheight=15pt]{geometry}
\usepackage[T1]{fontenc}\usepackage{lmodern,amsmath,amssymb,fancyhdr,parskip,hyperref}
\hypersetup{colorlinks=true,urlcolor=blue}
\pdfinfoomitdate=1\pdftrailerid{}\pdfsuppressptexinfo=15
\pagestyle{fancy}\fancyhf{}\lhead{Week 79 / Numbered-ball detectives}\rhead{Facilitator draft}\cfoot{Bellingham Math Circle / Week 79 / facilitator / \thepage}
\setlength{\parskip}{7pt}\setlength{\parindent}{0pt}
\begin{document}
\section*{Mathematical overview}
Under the smallest-label rule, after completed day $n$, $A_n=\{n+1,\ldots,2n\}$ and $B_n=\{1,\ldots,n\}$. Both counts are $n$. Every fixed positive label $k$ is introduced on day $\lceil k/2\rceil$ and is in $B$ permanently by day $k$. Under the \emph{eventual-membership} interpretation, the limiting sets are empty $A$ and all positive labels in $B$. There is no last finite day.

If instead we transfer the larger of the two new labels, after day $n$ the odd labels through $2n-1$ remain in $A$ and the even labels through $2n$ are in $B$. The finite counts still tie. Count histories do not determine limiting membership.

Experiments establish small cases. Induction and a named label's finishing day explain all cases. Children are meant to distinguish ``each label eventually'' from ``all labels by one day.'' Grades 2--3 can enact and follow named cards; Grades 4--5 can formulate the uniform-day obstruction and changed-rule comparison. This is a new teaching adaptation, not a claim that actual infinitely many days have occurred.
\newpage
\section*{Preparation and a first route}
For each pair: cards 1--24 and either two shallow bowls OR the supplied A/B mat as the single shared state; place unused cards in consecutive pairs outside both bowls. One adult anchors each participating table. For the three-child upper table, rotate feeder, mover and label detective. Adults record a few common observations; do not require children to track counts and every label at once.

Reading: short rules read aloud if needed. Arithmetic: compare labels, count to 24, add two and remove one. Reasoning: track one object across several changing arrangements. The limit discussion is optional; do not mistake arithmetic speed for readiness.

After the usual five-minute movement break, allow two minutes to handle/sort cards. Launch together for three minutes: demonstrate adding a pair, locating the \emph{smallest label anywhere in A}, and moving it to B. The printed before/intermediate/after visual supports the convention. Show a legal finite move, not the limiting answer.

First route: 15 minutes enacting the first six days and predicting arrangements; twelve days is an optional continuation; 10 minutes with children adopting labels; 10 minutes on the changed rule; five minutes of shared explanations. Leave at least five minutes for questions or replay. Use one mat repeatedly; bring in later questions when the table is ready. Draft timing and physical handling are untested.
\newpage
\section*{Facilitator reasoning and hints}
Induction: assuming $A_n=\{n+1,\ldots,2n\}$, adding $2n+1,2n+2$ and transferring $n+1$ gives $A_{n+1}=\{n+2,\ldots,2n+2\}$. The transferred label appends to $B$. The empty day 0 is the base case. Problem 1: adding two and moving one increases each bowl's count by one, so both have $n$ cards after day $n$. Day 6 gives $A=\{7,\ldots,12\}$ and $B=\{1,\ldots,6\}$. Day 10 gives $A=\{11,\ldots,20\}$ and $B=\{1,\ldots,10\}$.

Problem 2: label 3 enters on day 2 and moves on day 3; 6 enters on day 3 and moves on day 6; 10 enters on day 5 and moves on day 10; 20 enters on day 10 and moves on day 20. For any $k$, entry is $\lceil k/2\rceil$ and transfer is day $k$. When a child adopts label $k$, ask only when that card enters and when it leaves A. For example label 8 enters on day 4, remains in A through day 7 and moves on day 8. An unseen label can be discussed with a placeholder; the limited card supply does not stop the mathematical rule.

If the table says A is always equally full, agree that the finite count claim is correct. Ask what happens to one \emph{fixed} label. Problems 3--4: inspect a named finite list of labels, leaving every A/B card where it is. Inspect labels 1--6, then 1--12, regardless of location. Any fixed finite list eventually has all its labels in B. At every positive finite day $N$, label $N+1$ is still in A. Day 0 is empty because no labels have entered, and is not a day when all positive labels are in B. Hence no finite day finishes every label.

For the changed rule, the larger new card is $2n$, so it moves immediately to B. The smaller $2n-1$ never moves. Have children use the same finite number of days and compare the two actual sets. Do not print the odd/even answer on their working board.
\newpage
\section*{Limits, encore, sources and observation record}
Adult analysis: the indicators $a_n(k)=1_{A_n}(k)$ tend pointwise to 0, while $\sum_k a_n(k)=n\to\infty$. Thus $\sum_k\lim_n a_n(k)=0$ differs from $\lim_n\sum_k a_n(k)$. Convergence is not uniform: $\sup_k a_n(k)=1$ for every $n\ge1$ (day 0 is empty). Infinite counting measure and the lack of a summable majorant matter; nonuniformity alone is not a universal explanation for every failed interchange. For B, membership is monotone and its count tends to infinity. None of this establishes an extra ``midnight state'' without selecting the membership interpretation.

Encore: children invent a transfer rule while keeping the two-per-day/one-transfer convention. Investigate labels that eventually move, labels that never move, and labels whose locations never settle if backward transfers are newly permitted. State any changed rules explicitly. Connect later to the midpoint process in Week 80, or existing Week 50 staircase limits. Those are different questions about limits.

Source: the organizer's numbered-ball process and changed-rule proposal from the October 9 brief; finite formulas, comparison and proof are independently derived here. Pedagogy: Rozhkovskaya, \emph{Math Circles for Elementary School Students}, Lesson 8, ``At the lesson,'' item 2 (EPUB part0018): instructors reduced slow copying. Our adult-scribed record is an adaptation.

Unpiloted. Print US Letter, single-sided, Actual Size. Physical fit, card handling and classroom procedure have not been rehearsed. Record afterward: tasks actually used; child actions and words; confusion requiring rescue; questions they wanted to continue. Separate observations from explanations or hypotheses.
\end{document}
